\textbf{Решение.} Расположим действительную ось вдоль прямой AC так, чтобы точки $A$ и $C$ имели координаты $-1$ и $1$ соответственно. Тогда из условия $AB=BC$ следует, что точка $B$ лежит на серединном перпендикуляре к $AC$, то есть $B=ib$ для некоторого $b>0$. Поскольку $L$ лежит на диагонали $AC$, её  координата имеет вид $l\in(-1,1)$. Так как точки $B,L,D$ коллинеарны и $L$ лежит между $B$ и $D$, можно записать $D=l+u(l-ib)$, где $u>0$.

Пусть $x$ и $y$ --- координаты точек $X$ и $Y$. Так как $X\in AB$, существует вещественное $\lambda$ такое, что $x=-1+\lambda(1+ib)$. Кроме того, точки $B,C,D,X$ лежат на одной окружности, а потому отношение
\[
\frac{(x-ib)(1-d)}{(x-d)(1-ib)}
\]
вещественно, где $d=l+u(l-ib)$. Подстановка $x=-1+\lambda(1+ib)$ и упрощение дают
\[
(\lambda-1)\bigl((1+b^2)(1-l)\lambda-2(1-l^2-u(b^2+l^2))\bigr)=0.
\]
Корень $\lambda=1$ соответствует точке $B$, поэтому
\[
\lambda=\frac{2(1-l^2-u(b^2+l^2))}{(1+b^2)(1-l)}.
\]

Аналогично, поскольку $Y\in BC$, существует вещественное $\mu$ такое, что $y=ib+\mu(1-ib)$. Из того, что точки $A,B,D,Y$ лежат на одной окружности, следует вещественность отношения
\[
\frac{(y-ib)(-1-d)}{(y-d)(-1-ib)},
\]
и после подстановки и упрощения получаем
\[
\mu=1-\frac{2(1-l^2-u(b^2+l^2))}{(1+b^2)(1+l)}.
\]

Теперь введём обозначения $P=2u(b^2+l^2)+l^2-1$ и $Q=b(1-l^2)$. Тогда прямым вычислением из найденных формул для $\lambda$ и $\mu$ получаем
\[
x-l=-\frac{(1+ib)(P-iQ)}{(1+b^2)(1-l)},\qquad y-l=\frac{(1-ib)(P+iQ)}{(1+b^2)(1+l)}.
\]
Следовательно,
\[
(x-l)(y-l)=-\frac{(1+ib)(1-ib)(P-iQ)(P+iQ)}{(1+b^2)^2(1-l^2)}=-\frac{P^2+Q^2}{(1+b^2)(1-l^2)}<0.
\]
Итак, произведение $(x-l)(y-l)$ является отрицательным вещественным числом.

Так как точки $X$ и $Y$ лежат на сторонах $AB$ и $BC$ выпуклого четырёхугольника, они находятся в верхней полуплоскости, откуда $\arg(x-l),\arg(y-l)\in(0,\pi)$. Из условия $(x-l)(y-l)<0$ следует, что
\[
\arg(x-l)+\arg(y-l)=\pi.
\]
Но луч $LA$ направлен вдоль отрицательной вещественной полуоси, а луч $LC$ --- вдоль положительной, поэтому
\[
\angle XLA=\pi-\arg(x-l),\qquad \angle YLC=\arg(y-l).
\]
Значит,
\[
\angle XLA=\angle YLC.
\]
Что и требовалось доказать.